% P2 supplementary method: decision and numerical validity boundaries.
\subsection{A conservative, content-level safety rule and its limits}
\label{sec:p2-safe-grouping-rule}
Content grouping is performed \emph{before} creating
train/validation/test folds. For reproducibility, a candidate
epoch match is keyed by numeric dtype and byte order, array
shape and channel order, event label, time-index convention
and full numeric sample content. A cryptographic digest
accelerates the search; the matching raw canonical byte
payload is compared directly before claiming exact identity.
Two released IDs with any matching labelled epochs form
a cross-ID graph edge; connected components define a
conservative unit for \emph{content-disjoint} splitting.
This graph is a \emph{split hygiene constraint}, not a
claim that a component necessarily corresponds to a
single person, a single recording session, or a
single clinical recruitment identity. A partial
one-epoch collision need not establish identity
of full recordings.

We distinguish a complete epoch-set alias, as found
between the specified released IDs 4 and 5, from
isolated epoch overlap. The latter requires separate
investigation of stimulus alignment, repeated
calibration sequences, resampling and deterministic
synthetic trial insertion. Hash equality without
a full canonical-byte comparison, equality of
target AUC values, or similarity after arbitrary
rounding is \emph{not} our identity criterion.

\paragraph{Robustness of the reported aggregate effect.}
The matched 49-group experiment isolates a
\emph{pipeline} contrast, not a causal intervention.
For source-macro AUC, one affected group contributed
$-0.2625$ whereas the remaining 48 groups averaged
$-0.00078125$; thus the arithmetic overall
$-0.00612245$ is dominated by a single component.
Removing that group's contribution from the
reported average is an influence diagnostic,
not a new independent 48-group refit.
A new clean acquisition would be necessary to
estimate population-wide consequences of
record aliasing, and a clarification from
the public-data custodians is required
before making any attribution about the
original acquisition process.

A standalone synthetic split-safety regression
test checks exact-byte confirmation, label
sensitivity, full-trial aliasing and the
published factorial \emph{arithmetic}.
It does not read private patient identifiers,
reproduce the public EDF files, re-estimate
the published AUCs or establish new validation.
